% ============================================================
%  H. Høffding, "Vort Hjem. En Indledning"
%  Indledning til: Fru Emma Gad (red.),
%  Vort Hjem, Bind I.
%  København: Det Nordiske Forlag, Ernst Bojesen, 1903.
%  TRANSLATION (English) --- Hans Halvorson
%  [Sections I--II translated; remainder forthcoming]
% ============================================================
\documentclass[12pt,a4paper]{article}
\usepackage[utf8]{inputenc}
\usepackage[T1]{fontenc}
\usepackage{libertinus}
\usepackage{libertinust1math}
\usepackage[english]{babel}
\usepackage[protrusion=true,expansion=true]{microtype}
\usepackage{setspace}
\usepackage[margin=1.2in]{geometry}
\usepackage{fancyhdr}
\usepackage{hyperref}

\hypersetup{
  pdftitle={Our Home. An Introduction},
  pdfauthor={H. Høffding},
  hidelinks
}

\pagestyle{fancy}
\fancyhf{}
\rhead{Høffding, \emph{Our Home} (1903)}
\cfoot{\thepage}

\setstretch{1.3}

\title{\textbf{Our Home}\\[6pt]
  \large An Introduction}
\author{Harald Høffding\\[4pt]
  \small Translated from the Danish by Hans Halvorson}
\date{Introduction to: Mrs.\ Emma Gad (ed.),
  \emph{Vort Hjem}, vol.\ I\\
  Copenhagen: Det Nordiske Forlag, Ernst Bojesen, 1903}

\begin{document}
\maketitle
\thispagestyle{fancy}

\bigskip

\section*{I.}

\noindent The word \emph{home} calls up associations of intimacy,
fullness, and rest. Of intimacy --- for although the home, viewed
from outside, is a place like all other places, it is for us the place
where we received everything that fate brought us: the inlet to which
the effects of the great ocean of the outer world had to find their
way before we could feel and appreciate them. And it is equally the
place from which the effects we ourselves have exercised have gone
forth. In all our acting and suffering, all our giving and receiving,
the mood of home therefore weaves itself in and forms a background
of which we perhaps first become clearly conscious when it is
threatened with dissolution, or when it undergoes some deeper change.
And with this the fullness is also connected. For it is a large circle
of memories and habits --- indeed, the whole content of our life ---
that coheres with this central idea. Home is like the net that cannot
be drawn up without a swarm of living content following with it. Rest,
finally, depends on both intimacy and fullness. Rest cannot be found
in the foreign and the new, only in what is known and tried, and
whose appropriation requires no exertion. What is homely is fitted to
the forms of our nature; every fold and joint in our being finds there
a resting place that suits it. When what is homely surfaces again for
us, it is as though it had never been away; easily and quickly we
settle back into the long-since-fitted frame. And the wealth of habits
and memories that home bears also brings rest. One cannot rest on a
sharp edge; there must be a manifoldness in the content if continual
dwelling on one and the same thing is not to tire and dull.

One has a home as soon as one has something to lay one's head against.
Home presupposes fixedness, constancy, repetition. What is homely
stands in opposition to the new, the foreign, and the shifting. When
the fox has its lair and the bird its nest, they have a home. They
have something to return to, something that waits for them, something
that can be recognized and that stands firm through the surge of
change. We therefore also have a home in our own inner life when
the content of our consciousness does not shift without interruption,
but we have fundamental thoughts and fundamental moods that
constitute our spiritual standpoint and condition the effect that new
experiences have on us --- indeed, that are the very precondition for
anything being able to appear as new to us. Into the circle of these
fundamental thoughts and moods everything must be received that we
would genuinely call our spiritual possession.

It is more than an arbitrary comparison that we draw here when we
place the outer, material home and the inner, spiritual home alongside
each other. For they both rest on the same physiological and
psychological conditions. Life cannot persist without a rhythm in its
inner and outer relations. Intake and output, inhalation and
exhalation, rest and movement must alternate with one another if life
is to persist. A return to the same states is therefore an organic
condition of life. And spiritually, the course of life requires fixed
points of departure and orientation, stable background-states against
which experiences can stand out in their distinctiveness --- yet in
such a way that a common measure becomes possible.

\section*{II.}

\noindent In the human world, home and family belong closely together.
Yet there can be a home where there is no family. The solitary person
can have a home, and at the lowest levels of human life known to us,
the family has not yet separated itself from the horde or clan as a
distinct community. The wandering horde returns by preference to the
places that offer sustenance and security, and in any case carries its
tents and utensils along as constituents of a travelling home. Already
here home means intimacy, fullness, and rest. The struggle for life,
which is a struggle against the foreign, comes to a halt in the tent
and on the familiar pastures and hunting grounds. And memories
already play a role when one moves through the same regions where
one's ancestors moved. Even if one has no fixed dwelling places
otherwise, the graves of one's forebears will be a center around which
one circles and for which one takes up the fight, even when one offers
no resistance against pressing foreign forces elsewhere. Already here,
then, home undergoes the significant metamorphosis from being a
material necessity to becoming an ideal value. What from the beginning
was only a physical refuge can later become the only place where life
seems worth living, because it can there alone be lived in continuity
with the life of the past, as a continuation of it.

Of the greatest significance in the history of home was the development
of the family into a distinct kind of community. The core of the family
is the bond between mother and child. These two hold together and
have their home with each other, however things may stand with the
family's other elements. Here a natural connection is at work --- not
merely the connection conditioned by the satisfaction of the need for
sustenance and security in the same places. Birth is an extension of
organic growth beyond the individual organism. Home thus grows,
literally. In many parts of the world one finds, at a certain stage of
development, the arrangement by which husband and wife continue to
belong to their respective tribes, and the children belong to the
wife's tribe, not the husband's. What has created the family as we now
know it --- as a community of husband, wife, and children --- has been
partly the man's will to power and self-assertion, and partly the fixed
dwelling. The man regarded the woman and the children she had borne
him as his property; he wished to make use of the labor they possessed,
and he felt a need to see his lineage continued in his offspring. The
fixed dwelling and the regular shared life it brought about provided a
foundation for custom and usage, for tradition and memory, and there
was now more occasion for the arising of tenderer feelings. Shared life
softened, underpinned, and humanized the animal impulse and the crude
will to power. The common hearth became sacred, for the spirits of the
ancestors visited it continually. And in this sanctuary even the
foreign guest --- even if he were an enemy --- found a place of refuge.
A sanctuary had been created whose spirit could spread to wider circles.

The family, at its most primitive level, is only a means for the
continuation and perpetuation of life. But even here, interests and
motives of a different kind from those that originally founded the
bond begin to form. It is then not only sustenance, security, and the
continuation of the lineage that are sought; the mood of home that the
fixed dwelling and the shared traditions engender becomes valuable in
itself, stands as an immediate good in which the human being wishes
to live. A common fate and common labor continually consolidate the
feeling of belonging together. Even if the natural bond only comes
fully into view in the relation between mother and child, bonds are
woven between all the members of the family in the course of their
shared life that can often rival the merely physical connection in
strength.

Even after the family has separated itself from the horde or clan and
has become a community in its own right, sustained by its own
distinctive forces, it does not yet stand in the opposition to other
kinds of community that we know from modern life. It has more
purposes than the modern family, since it is at the same time a small
state and a small church, a legal community and a cultural community
in one. Paternal and parental authority holds sway as the highest
power, and even where a state exists, the family threshold forms a
boundary to its influence. Religion, which at this stage is the sum of
all culture, is from the beginning the affair of the family: within the
family the spirits of the deceased are venerated from generation to
generation; the hearth-fire is kept burning in their honor, and their
commands constitute the content of the moral law. Only under the
advance of development does the family cease to have such wide-ranging
purposes; the state draws all relations among individuals under its
power and supervision, and distinct communities form for the
cultivation of material and ideal culture. What then remains for the
family is the tending of the first shoots of new human life, and of the
intimate feelings that can bind human beings together apart from all
cultural endeavors and all legal relations. This in no way diminishes
the family's significance. On the contrary: the family relation, by
virtue of its independence from all the differences and antagonisms
that in wider circles divide human beings, and by virtue of its
connection with the natural instincts that secure the continuation of
the lineage, stands as the strongest and most intimate community ---
one that provides a measure for the development of all other
communities. The more intimate and stronger other social bonds become,
the more they are described in terms borrowed from the family's. In
the family the individual can live according to his whole personality,
whereas the cultural community and the state have use only for certain
sides of his being. And the bond here is grounded in the unconscious
natural life, whereas other communities rest on half- or wholly
conscious purposes and interests.

Even where the division of labor brings about an opposition between
the family on the one side and the cultural community and state on
the other, the family therefore retains its significance --- or indeed
first comes fully into its own. It now becomes a central place where
what is accomplished by work in the cultural community and state can
be appreciated and enjoyed. And it becomes a nursery where the forces
that are to work in the cultural community and state can germinate
and grow. It represents intimacy, coherence, security, and rest, to
which one can return again and again from the outward-turned, divided,
struggling, and strenuous activity in the wider circles.

\section*{III.}

\noindent We have described home according to its ideal as well as
according to its historical development. But it emerged that a great
many different forces and conditions are presupposed if the ideal is
to be reached here, and it is therefore to be expected that the ideal
is not always reached. And if an old saying is right --- that the
better something is, the worse it becomes when corrupted --- then we
must expect that home under unfavorable conditions will oppress life
just as much as it can advance life under fortunate ones.

Let us first examine somewhat more closely what the enduring
significance and justification of home rests upon. The justification
lies in a line of thought similar to the one that justifies the right
and freedom of the individual personality. For the sake of society
itself, the greatest possible independence and freedom of the
individual must be secured. For in the individual's inner life lies
the seed of everything new and great. The more vigorously life stirs
there, the more vigorously it will be able to stir in the wider
circles. Individual personalities are the source-points of society's
life, just as they are the central places to which all effects
originating from social life propagate themselves. There the value of
life is felt, and from there proceed the endeavors that can enhance
that value. Smaller and narrower communities stand in a similar
relation to larger ones. Life is conducted and felt more strongly and
intimately in the former than in the latter. What is not attended to,
and what is not protected, in the wider circles finds attention and
care in the smaller ones. For this reason the interests of the large
communities --- those of the state and the cultural community --- themselves
require that the smaller ones be maintained in their distinctiveness
and independence. The significance of home is likewise grounded in
the inverse relation between the strength and the scope of life.
Already Aristotle, in his critique of Plato's \emph{Republic} ---
where the family for the higher classes was to be abolished in favor
of a universal kinship among all individuals --- appealed to this law.
It will, he says, become a diluted love that is supposed to unite
all after the dissolution of the narrower and more intimate bonds;
it will scarcely be perceptible, any more than the admixture of a
little sugar in a large quantity of water makes an impression on the
taste. It is perfectly correct, Aristotle continues, that the state
should be a unity; but if one erases all difference and multiplicity,
the state ceases to exist, just as the symphony ceases when only one
voice is heard, and the rhythm when there is only a single verse-foot.
--- The great Greek philosopher of the state thereby clearly
demonstrated the significance of the small communities' existence
within the large.

But in order to understand the doubts that may arise about the value
of home, we must take note that the law of the inverse relation
between strength and scope is not the only law governing the
development of human interests. For in the domain of emotional life
there also appears an inverse relation between strength and content.
The content of a feeling is the representations to which it is
attached, and which give it its distinctive character. Joy and sorrow
are nuanced and varied according to the representational content that
conditions them. Now it is clear that with narrow scope there also
easily follows a narrow content. In the small circles, emotional life
easily assumes a purely elementary character, coming to lack definite
and many-sided formation. When isolated from life in the wider
circles, there then easily enters --- even if the strength of the
feeling is preserved, which on account of the influence of habit and
monotony does not always happen --- a spiritual poverty. The feeling
then becomes blind. The content it has, which is perhaps held fast
with great intimacy and passion, derives essentially from the past.
The household gods of the family are the heroes of bygone ages; but
it is not given that they are suited to be guides in new times. The
conservative character of home is conditioned by this. Intimacy,
security, and rest then turn into inertia and engender reluctance and
fear in the face of the new. What otherwise constitutes the strength
of home here becomes its weakness. All ages of history offer
experiences in this direction. Here we will draw attention only to a
few particularly striking features. --- A Japanese writer, Tokiwo
Yokoi, recently undertook, on the occasion of the war between China
and Japan, a comparison between the prevailing morality of the two
peoples, and laid great stress on the fact that Chinese morality is
essentially a family morality: ``It was the principle of filial piety
that formed the foundation of Chinese morality, rather than the
principle of loyalty to the sovereign.'' The Chinese ethical tradition
was pushed all the more in this direction because dynasties in China
have changed very frequently. The family bond then became the only
bond that could uphold unity and continuity in society --- and with
such firmness that ``these three hundred million people have preserved
their traditions and customs unchanged for three thousand years and
through twenty-three changes of dynasty.'' But the virtues that
distinguish the Chinese are therefore also only the virtues of
private and domestic life, whereas patriotism, and integrity in public
affairs, are lacking. This explains the difficulty of taking up and
carrying through reforms in China, while the Japanese, who through
long ages had been practiced in subordinating private and domestic
interests to the demands of the state, threw themselves onto the path
of progress with the greatest zeal and self-sacrifice. --- Discerning
observers of conditions in Turkey and Persia have found that the
oriental harems perhaps offer the most important resistance to the
penetration of a new culture. ``While there,'' says Vambéry
(\emph{Der Islam des neunzehnten Jahrhunderts}), ``in the Selamlik
(the part of the house accessible to men and the outside world) signs
of renewal appear here and there, and \emph{\`{a} la franca} has so
to speak become a watchword, in the harem, the sanctuary of the
Muhammadan family, everything has remained as it was, and not the
slightest change has occurred in the furnishings, customs, and
household order, any more than in dress, manner of speech, and
habits of thought\ldots\ Among some hardened Orthodox, the harem
has in reality become a happy refuge, while it is of course a
tremendous impediment to those champions of the new who are in
earnest. What is introduced in the Selamlik is either ignored in the
harem or torn down again by unbridled fanaticism.'' --- Tendencies in
a similar direction are not lacking in Europe. One has thus ---
both in Protestant and in Catholic countries --- advanced the claim
that the school should be dependent on the home. ``Home, state, and
church,'' said Bishop Ketteler, ``must have schools that correspond to
their spirit and their demands.'' What is missing here is the
supplementary proposition that the school --- in the widest sense:
the community for the advancement and appropriation of science ---
must in turn act back upon the spirit and the demands that prevail in
home, state, and church; that there is a wider horizon that the home's
(and the other communities') tendency toward closure easily leads it
to restrict. The often close and stale air of home needs to be
cleansed by fresh currents. In such cases one is not in fact fighting
for the home, but using the home as a means of preserving dogmas
that cannot be maintained in any other way. But the cultural content
on which the home is to draw will by this procedure become poorer
and will congeal in dead forms.

It will in the long run not be possible to preserve strength and
intimacy when the content shrinks in this way. Repetition will
exercise its deadening power, and spiritual death will be the final
effect. All living feeling presupposes for its persistence rhythm and
opposition; where there is nothing new to discover at all, feeling
passes over into passive habit. The character and value of home will
therefore depend on the richness in the nature of the people who
constitute it. The more distinctive and many-sided each of them is,
the more favorable will be the conditions for new discoveries within
the narrow circle, and the more independent the home will be able to
become of the wide world, without losing life and freshness.

But it is precisely not always the case that the individual personality
finds in the home the right conditions for its development, or meets
with the understanding without which its distinctiveness cannot unfold.
In the home it is naturally the common that holds sway. The natural
ground that sustains the family holds the individual individualities
fixed in their still-undetermined form; in the more determinate
formation of the individual individuality there lies a possibility that
the immediate connection and the instinctive mutual understanding
will cease. A reluctance to break with the shared life naturally arises
--- to turn aside from the customary tracks along which the familiar,
unwritten-law-governed outlook guides the members of the family ---
and the individual is easily impeded thereby in his independent
development. In its most prosaic form, this reluctance appears as
shyness. A peculiar kind of self-consciousness that comes into force
over against those nearest, but can fall away over against those more
distantly placed. This relation forms an analogy to the fact that the
feeling of love, at higher levels of human life, does not awaken
within the sibling relationship but requires a new and foreign object
to be aroused. Only in the face of the new and foreign does
individuality come fully forth. Add to this that within the circle of
home both the unfortunate and the fortunate elements of character can
thrive. What is easily worn smooth in contact with the less
indulgent foreign world can, under the mild and soft conditions of
home, spread and draw attention away from what is essential and
valuable. In any case, within the home the petty and the great often
emerge in obscure mixtures that make both overestimation and
underestimation possible. Now the individual is regarded without
justification as a prophet by those nearest to him; now he is not
respected as a prophet, though he is one.

\section*{IV.}

\noindent The future of home will depend on whether the dangers
pointed out here can be avoided. If home is to uphold its value for
human life, it must enter into an open and positive relation to the
wider communities and the interests stirring in them, and it must
also within its own circle bring about the most favorable possible
conditions for the distinctive development of its individual members.
Home can be attacked from two sides: from the socialist side as
isolating, from the anarchist side as inhibiting. Human development
moves at once in the direction of wider horizons and more universal
spheres of activity, and in the direction of a more determinate
formation of individual distinctiveness.

But there is no reason to think that home will not be able to meet
the tasks set before it. It has hitherto followed the human being
along the long road from animal rawness to social and spiritual human
life, and even if the household gods change --- and may perhaps change
more frequently in the future than before --- their place will not come
to stand empty.

It is a natural necessity that we live in a small world before the
large one can become real for us. And the more the home truly is
itself a small world, the more can life there become a preparatory
school for life in the wide world --- while at the same time
preserving its independent value as the place of intimacy and
security. There is something common to every world, be it small or
large. A common law holds sway and determines the place and activity
of the individual elements; each individual element must subordinate
itself to the order of the whole, yet is also animated by the bonding
forces that hold the whole together. Within the family, the individual
unconsciously grows into the idea that he is one among many; he comes
to know self-surrender, resignation, and justice; and without noise,
without any painful resistance needing to be felt, the unbridled
expression of his egotism is restrained. The first objects of fear and
love, of admiration and gratitude, of esteem and reverence, he finds
within the circle of home.

Even within the small world --- when it does not isolate itself ---
the individual will have experiences of the significance of larger
relations. The swell from the ocean of the world often carries into
the quiet inlet of home. National misfortunes, for example, touch
the child early through the sorrow they awaken in the elders; to its
amazement it sees that what happens in the distance affects those
nearest, as though it were their own weal and woe at stake. The first
feeling of belonging to a larger whole than that constituted by
parents and siblings, family and friends, is perhaps laid down in
such a moment. Through its relation to the first objects of its love
and admiration, the child unconsciously grows into the interests for
which these, its models, already live. But the elders stand not only
as intermediaries between the child and the wider circles of interest.
They also stand as intermediaries between the past and the future.
Their task is not merely to transmit the values of the past, but also
to prepare those of the future. Home has here the great significance
of bridging the different generations and making possible an
understanding between them that would be excluded if the new
generation did not grow up within the older one and under the
influence of the most immediate and most naturally determined
sympathy that can in general bind human beings together. The
struggle between the old and the new would be sharper and also less
fruitful if it were not prepared in subdued form within the walls of
home.

There are also within home the most favorable conditions for personal
development, especially in the period when the forms are still
undetermined and many shoots are putting forth that are perhaps not
viable, yet cannot be forcibly cut away without danger to life. There
further lies in the nature of sympathy, as it can appear within home,
an interest in and a capacity for understanding purely individual
stirrings, which is excluded only under unfavorable circumstances,
when blindness and limitation come to hold sway. In any case,
individual distinctiveness will scarcely find such favorable conditions
for life as here. Greater demands will also in this respect be placed
on the home of the future. It is after all now --- to bring forward a
single important point --- not only the man whose personality demands
free scope for individual distinctiveness. The same demand is now
being raised for women as well. Women now wish not only, as could
also happen in earlier ages --- in the Renaissance period, for
instance --- a share in the general education and in the thought of
the age, but also the possibility of a vocation other than that
exercised in the domestic domain itself. The great law of the division
of labor will also make itself felt for women's education and women's
work. While woman within the home --- as mother, wife, and sister ---
could be said to represent what is undividedly human, the enduring
wellsprings of intimacy and love, the law of the division of labor
and one-sidedness now threatens them as well. Many see in this a
great danger for the future of home. But if home is to be capable of
solving the task of allowing personal distinctiveness to move more
freely than hitherto for the man's part, the corresponding task will
also be capable of being solved for the woman's part. Home has so
deep and firm a natural foundation that it will certainly endure, even
if the hitherto prevailing division of labor between the two sexes is
replaced by a different arrangement. Nature has not yet spoken its
last word here.

Like everything else in the world, home too must struggle for life.
Under very different forms it has hitherto maintained its significance,
and it will also in the future be able to demonstrate its great
significance: to be at once a world of intimacy and a world of freedom
to a far greater degree than any other kind of community is able.

\end{document}
%%% Local Variables:
%%% mode: latex
%%% TeX-master: t
%%% End:
